\bookchapter{פרק 61: Swift למתחילים — השפה הראשונה לפיתוח באפל}

\chapterintro{בפרק הזה נתחיל סוף סוף ללמוד את השפה. אל תילחצו — Swift ידועה בידידותיות שלה למתחילים, ובנוסף אין צורך לשנן ספר תחביר: שליטה ב-20\% מהתחביר מספיקה כדי להבין ולכתוב 80\% מהקוד היומיומי. את השאר תחפשו תוך כדי עבודה על הפרויקט. זכרו את עקרון הלמידה שלנו: להבין עדיף על לשנן, לדעת להשתמש עדיף על לשלוט לעומק.}

\booksection{משתנים וטיפוסים: קוראים לנתונים בשם}

תוכנה היא “נתונים + לוגיקה”. משתנה הוא בעצם לתת לנתונים שם ולאחסן אותם. ב-Swift מצהירים עם let ו-var:

\begin{lstlisting}[language=Swift]
let name = "דני"        // let: קבוע (מומלץ כברירת מחדל)
var age = 18             // var: משתנה (אפשר לשנותו אחר כך)

age = 19                 // <i class="fas fa-check-circle mr-1"></i> var אפשר לשנות
// name = "נועה"         // <i class="fas fa-times-circle mr-1"></i> let לא ניתן לשינוי, שגיאת קומפילציה

let score: Int = 100     // אפשר גם לציין את הטיפוס במפורש
let pi = 3.14            // גם בלי לכתוב אותו הוא יוסק אוטומטית

print("שלום \(name)")   // שילוב מחרוזות: לוכסן אחורי + סוגריים
\end{lstlisting}

את הטיפוסים הנפוצים אפשר לזכור בטבלה אחת:

\begin{center}
\begin{tabularx}{\linewidth}{|X|X|X|}
\hline
\textbf{טיפוס} & \textbf{מה מאוחסן} & \textbf{דוגמה} \tabularnewline
\hline
Int & מספר שלם & 42 \tabularnewline
Double & מספר עשרוני & 3.14 \tabularnewline
String & טקסט & "שלום" \tabularnewline
Bool & אמת/שקר & true / false \tabularnewline
{}[Int] (מערך) & קבוצת נתונים & {}[1, 2, 3] \tabularnewline
\hline
\end{tabularx}
\end{center}

\begin{center}\small\color{gray}הטיפוסים הבסיסיים הנפוצים ב-Swift\end{center}

\begin{tcbinfo}{טיפ: כלל העדיפות של let}
חוק הברזל של קהילת Swift: ברירת המחדל היא let, ו-var רק כשזה באמת נחוץ. אם תשתמשו ב-var במקום שאפשר let, Xcode יזהיר אתכם בטקסט צהוב. ההרגל הזה חוסך לכם שלמה שלמה של באגים מהסוג “משתנה ששונה בטעות” — כמה שפחות דברים משתנים בקוד, כך קל יותר להבין אותו.

\end{tcbinfo}

\booksection{ערכים אופציונליים (Optional): “הקמיע נגד התרסקות” של Swift}

זה המושג החשוב ביותר ב-Swift, בלי תחרות. במשפט אחד: כשערך עלול “לא להתקיים”, Swift מחייבת אתכם לטפל בו — השפה מחסלת מהיסוד את סיוט התעשייה של “התרסקות בגלל מצביע ריק”.

\begin{lstlisting}[language=Swift]
var name: String = "דני"   // טיפוס רגיל: לעולם לא יכול להיות nil
// name = nil                // <i class="fas fa-times-circle mr-1"></i> שגיאת קומפילציה: עוצרת אותך מיד

var nickname: String? = nil  // מוסיף סימן שאלה: מאפשר nil (ערך optional)
nickname = "טומי"

// עם ערכי optional חובה "לפתוח" אותם:
if let n = nickname {        // if let: פתיחה בטוחה
    print("שלום \(n)")
} else {
    print("עוד אין כינוי")
}

let len = nickname?.count ?? 0   // ?? ערך ברירת מחדל (אופרטור coalescing)
print(len)
\end{lstlisting}

\begin{tcbwarn}{שימו לב: ! היא "פעולה מסוכנת"}
לכתוב בקוד nickname! (פתיחה כפויה) זה כמו לצעוק למהדר “אני מבטיח שזה לא nil” — וברגע שטועים, האפליקציה מתרסקת על המקום. כלל למתחילים: השתמשו כמה שאפשר ב-if let / ?? ואל תשתמשו ב-!. לפעמים Xcode מוסיף ! מעצמו בעת השלמת קוד — זכרו לתקן את זה.

\end{tcbwarn}

\booksection{פונקציות וסגרים: אורזים לוגיקה ב“כפתורים”}

פונקציה היא קטע לוגיקה עם שם, כמו כפתור בשלט רחוק. סגר (Closure) הוא “פונקציה בלי שם”, שאפשר להעביר לאחרים כמו שהוא:

\begin{lstlisting}[language=Swift]
// פונקציה: השם בקדמת; פרמטרים וטיפוס ההחזרה מאחור
func greet(name: String) -> String {
    return "שלום, \(name)!"
}
print(greet(name: "דני"))     // בקריאה, עם שם הפרמטר

// closure: מוגדרת במקום, קלילה וישירה
let numbers = [1, 2, 3, 4]
let doubled = numbers.map { $0 * 2 }   // $0 = האיבר הנוכחי
// [2, 4, 6, 8]

let big = numbers.filter { $0 > 2 }    // סינון
// [3, 4]
\end{lstlisting}

ב-SwiftUI סגרים בכל מקום: הסוגריים המסולסלים אחרי Button(“לחצו עליי”) \{ ... \} של כפתור הם סגר. שננו את ה“קומבו” של האוספים — map, filter, sorted — כי 90\% מעיבודי הרשימות בעתיד יהיו שילובים שלהם.

\booksection{מבנים ופרוטוקולים: שיטת “הלגו” של Swift}

ההעדפה הרשמית של Swift ברורה: מעדיפים מבנה (struct), ומחלקה (class) רק כשנדרשת ירושה. מבנה הוא “טיפוס ערך” — כשמעתיקים אותו כל עותק חי את חייו שלו, מה שמפחית מטבעו סוג שלם של באגים:

\begin{lstlisting}[language=Swift]
// ה-struct: קופסת הנתונים
struct Record {
    var id = UUID()
    var amount: Double
    var category: String
    var date: Date

    // ל-struct יש "memberwise initializer" מובנה: אפשר ליצור ישירות
}
let r = Record(amount: 25.5, category: "אוכל", date: Date())

// פרוטוקול: הסכם על "אילו יכולות חייבות להיות"
protocol Identifiable {
    var id: UUID { get }
}
struct Record: Identifiable { ... }   // לאמץ את הפרוטוקול = להתחייב למלא את דרישותיו
\end{lstlisting}

ב-SwiftUI כמעט הכל הוא struct (View הוא פרוטוקול). אם המשפט הזה לא ברור כרגע — לא נורא, נפגש איתו יום-יום בפרק 62; בעולם של Swift מתחבקים קודם כל עם struct.

\booksection{כרטיס תחביר מהיר בשביל מי שמתחיל מאפס}

\begin{center}
\begin{tabularx}{\linewidth}{|X|X|}
\hline
\textbf{מה רוצים לעשות} & \textbf{כיצד כותבים} \tabularnewline
\hline
הגדרת קבוע/משתנה & \texttt{let x = 1}/\texttt{var x = 1} \tabularnewline
שיבוץ ערכים במחרוזת & "שלום, \textbackslash{}(name)" \tabularnewline
הגדרת פונקציה & func שם(פרמטר: טיפוס) -> טיפוס החזרה \{ ... \} \tabularnewline
תנאי & if תנאי \{ ... \} else \{ ... \} \tabularnewline
מעבר על איברים & for x in מערך \{ ... \} \tabularnewline
פתיחה בטוחה & if let x = ערך אופציונלי \{ ... \} \tabularnewline
ערך ברירת מחדל & x ?? ערך ברירת מחדל \tabularnewline
הגדרת מבנה & \texttt{struct X \{ var a: Int; var b: String \}} \tabularnewline
סגר & \texttt{\{ \$0 * 2 \}} \tabularnewline
\hline
\end{tabularx}
\end{center}

\begin{center}\small\color{gray}כרטיס תחביר מהיר ל-Swift (מהדורת מתחילים)\end{center}

\begin{tcbpractice}{הידיים על ההגה: 15 דקות של משחק עם Swift ב-Playgrounds}
יש מק? פתחו Xcode Playground. אין מק? השתמשו ב-Swift Playgrounds ב-iPad או במגרש המשחקים בדפדפן. הריצו כל קטע קוד בפרק, שחקו עם המספרים והטקסטים, וגרמו בכוונה לכמה שגיאות כדי לראות איך Xcode מתריע. 15 הדקות האלה של “לשחק לפני שלומדים” יקלו עליכם פי שניים כשתתחילו את פרק 62.

\end{tcbpractice}

\begin{tcbsummary}{סיכום הפרק}
\begin{itemize}
  \item מעדיפים let ונזהרים עם var; הטיפוסים מוסקים אוטומטית, ואפשר גם להצהיר במפורש אחרי הנקודתיים
  \item ערכים אופציונליים הם הנשמה של Swift: מה שעלול להיות ריק חייב טיפול, if let / ?? הם הקמיע ו-! היא פעולה מסוכנת
  \item בקריאה לפונקציה מציינים שמות פרמטרים; סגר = פונקציה בלי שם, והוא בכל מקום ב-SwiftUI
  \item מעדיפים struct ומקפידים על פרוטוקולים (Protocol); ב-SwiftUI כמעט הכל הוא struct
  \item 20\% מהתחביר עושים 80\% מהעבודה: קודם לדעת להשתמש, ואחר כך להעמיק
\end{itemize}

\end{tcbsummary}
